\documentclass[11pt]{article}
\usepackage{fontspec,amsmath,amssymb,amsthm,hyperref}
\usepackage[margin=1in]{geometry}
\setlength{\emergencystretch}{3em}
\title{Jensen repairs the averaged pairwise energy}
\author{LangevinLean mathematical exposition}
\date{}
\begin{document}\maketitle
\noindent Source module: \texttt{ProximalLangevin/ParticleEnergy.lean}. Compile this proof document with XeLaTeX or LuaLaTeX. The accompanying HTML places the complete Lean declarations beside every proof.


The scaled state is \(z=(\theta,x_1/\sqrt N,\ldots,x_N/\sqrt N)\), with \(N>0\), and \(J_i z=(\theta,x_i)\). The actual algorithm averages proximal parameter outputs; this module proves its energy inequality directly.


\section{Averaging squares in the shared parameter}


Cauchy–Schwarz for finite sums gives \(N(N^{-1}\sum_i b_i)^2\le\sum_i b_i^2\). Apply it separately to each parameter coordinate of the pairwise outputs. Each latent block instead satisfies an exact equality after multiplying by \(N\), since \((\sqrt N)^2=N\). Summing all coordinates yields

\[\|\operatorname{Pair}_P(z)-a\|^2\le\frac1N\sum_i\|P(J_i z)-J_i a\|^2.\]

\noindent Lean declarations: \texttt{average\_sq\_le}, \texttt{pairwiseMap\_sq\_dist\_le}.

\section{Sum genuine proximal energy inequalities}


The roadmap on v3 page 19 identifies averaged pairwise proximity with the proximal map of the lifted averaged potential. Shared parameter coordinates prevent that identity in general. The needed energy conclusion can nevertheless be proved without changing the algorithm.



For arbitrary pair outputs, Jensen's inequality for their averaged parameter gives


\[
\|\operatorname{Pair}_P(z)-a\|^2
\le\frac1N\sum_{i=1}^N\|P(J_i z)-J_i a\|^2.
\]


The scaled latent-coordinate contribution is an exact identity; only the parameter contribution requires an inequality. Apply the single-pair proximal energy result to every summand. Using \(\sum_i\|J_i(z-a)\|^2=N\|z-a\|^2\) and \(G_2=N^{-1}\sum_i g_2\circ J_i\) gives


\[
\|\operatorname{Pair}_{P_\lambda}(z)-a\|^2
\le\|z-a\|^2-2\lambda\bigl(G_2(z)-G_2(a)\bigr)+\lambda^2C.
\]


This lifted energy inequality is fully proved for actual minimizers. It repairs the deterministic proof bridge needed in Corollary A.13 while preserving the paper's particle algorithm and its original nonsmooth energy. It is still only one part of the full stochastic convergence proof.



The failure of the lifted-proximal identity has a concrete Lipschitz witness. Take \(N=2\), \(g_2(\theta,x)=|\theta-x|\), smoothing one, common input parameter \(1\) and physical particles \((0,4)\). The genuine pairwise proximal points are \((1/2,1/2)\) and \((2,3)\). Averaging only their parameter components gives \((5/4,1/2,3)\). In the lifted proximal objective with its scaled particle norm, the feasible point \((1,1,3)\) has value \(3/2\), while this pairwise output has value \(51/32\). The strict inequality is kernel checked below. The averaged energy proof above repairs the necessary estimate without this false identity.



\noindent Lean declarations: \texttt{pairwise\_prox\_energy\_le}.

\end{document}
